
Cross-architecture property prediction from neural network weights remains limited. Within-architecture methods achieve Spearman correlation $\rho$ = 0.81 for predicting ImageNet accuracy when transferring knowledge within the same architecture family (Schürholt et al., 2024). However, cross-architecture transfer---predicting properties of Vision Transformer models from knowledge learned on ResNet models---degrades to $\rho$ = 0.54, representing a 33\% performance drop \citep{Andreis2023}. This gap constrains practical applications where heterogeneous model zoos contain diverse architectures including ResNet, ViT, MobileNet, and EfficientNet variants.

The limitation stems from incompatible weight structures. ResNet employs convolutional tensors with spatial locality (C\_out $\times$ C\_in $\times$ H $\times$ W), while ViT uses attention matrices with global dependencies (d\_model $\times$ d\_model). Existing methods chose between two constrained solutions: architecture-specific encoding that preserves structure but requires shared token sizes across architectures (Schürholt et al., 2024), or architecture-invariant encoding that achieves cross-family transfer but discards operation-specific signals through uniform aggregation \citep{Andreis2023}.

This work proposes CAPE (Cross-Architecture Parameterized Encoder), which decomposes architectural differences into three learnable signals. Operation-specific encoders handle convolutional, attention, and MLP weights through modular mechanisms adapted from prior work: SANE tokenization for convolutions (Schürholt et al., 2024), UNF equivariance for attention \citep{Zhou2024}, and standard aggregation for MLPs. Contrastive projection aligns these heterogeneous embeddings through InfoNCE loss, leveraging the observation that models trained on the same task share objective functions despite architectural differences. Architecture graph neural networks encode computational topology---sequential connections with skip pathways versus dense intra-layer attention---as a learnable residual.

In proof-of-concept validation with 100 models and synthetic accuracy labels for mechanism validation, CAPE achieved $\rho$ = 0.67 on $ResNet\rightarrow ViT$ transfer, exceeding the target threshold ($\rho \geq$ 0.65) and improving 24\% over the SNE baseline ($\Delta\rho$ = 0.13, p = 0.032). All three components contributed independently without degrading each other. Diagnostic metrics confirmed that operation encoders produced distinct representations, contrastive alignment preserved architectural structure, and graph topology contributed meaningfully to predictions.

The results validate that architectural differences are not insurmountable barriers but separable signals amenable to modular learning. The proof-of-concept establishes mechanism function; full-scale validation with 400 models across four architectures and real ImageNet accuracy labels is required to demonstrate predictive capability.
